% 2026-11-23: Member D draft for the LNCS team report.
% Place the Method subsection in Method, and the remaining text plus the
% table and figures in Experiments and Results. Requires graphicx.
% Reference key moro2014bank:
%   Moro, S., Rita, P., Cortez, P.: Bank Marketing. UCI Machine Learning
%   Repository (2014). DOI: 10.24432/C5K306.
% Length: about one page with both figures. If the section budget is tight,
% drop the second figure (prior shift) and keep its numbers in the text.

% ------------------------------------------------------------------ Method
\subsection{Cost-Sensitive and Prior-Shift Evaluation}

Let $s_i$ be the log-odds assigned to record $i$ by a model, and let
$\hat y_i(t)=\mathbf{1}[s_i\ge t]$ be the decision induced by a cutoff $t$.
Ranking metrics such as ROC-AUC and trapezoidal PR-AUC summarise the whole
sweep over $t$, but every deployed rule commits to one $t$. We therefore
score an operating point by its expected cost,
\begin{equation}
C(t;r)=\mathrm{FP}(t)+r\,\mathrm{FN}(t)
\quad\text{per }1{,}000\text{ records},
\label{eq:d-cost}
\end{equation}
where $r$ is the cost of one missed subscriber relative to one wasted call.
Setting $r=1$ reproduces the symmetric error weighting that F1 and accuracy
imply. For each $r$ we report the cost of the best threshold per model, the
cost of the threshold that the validation split selected by maximising F1,
and the two trivial policies of calling nobody and calling everybody.

To probe the dependence on the base rate we reweight rather than resample.
Positives receive $w_{+}=\pi'/\pi$ and negatives $w_{-}=(1-\pi')/(1-\pi)$,
where $\pi$ is the observed positive rate and $\pi'$ the target one; weighted
precision, recall and both AUCs follow from the weighted cumulative counts.
No record is added, removed or relabelled, so the analysis isolates the
effect of the base rate on the metric itself. At $\pi'=\pi$ the weighted
metrics coincide with the unweighted ones to six decimal places.

% ------------------------------------------- Experiments and Results
\subsection{How Much of the Comparison Survives the Assumptions?}

We evaluate the five models of this report on the 6,781 test records
(793 subscriptions, 11.69\%) with the same split, feature encoding and
exclusion of call duration as the rest of the report, and no other change:
Gaussian NB on 50 encoded columns (member A), logistic regression in its
MLE and Gaussian-prior MAP forms on 41 columns with $\lambda=10$ (member B),
and the same NB on 45 PCA or SVD components (member C). Re-fitting all five
from their modules reproduces every published ROC-AUC and PR-AUC to within
$5\times10^{-7}$. Each cutoff is the one that maximises F1 on the 6,782
validation records.

\paragraph{Model choice is not the decisive variable.}
Under the cost of \eqref{eq:d-cost} the five models lie within
$0.65\%$ of one another at every ratio from $1$ to $20$: at $r=20$ the
cheapest model costs $789$ per $1{,}000$ clients and the most expensive
$839$, whereas the ranking between the two logistic fits alternates with
differences of a few tenths of a percent. The ordering is stable in one
respect only: the two logistic fits lead the naive-Bayes family at every
ratio, and the ordering of all five models is unchanged across the sixty
base rates examined below.

\paragraph{The cutoff is.}
Holding the models fixed and moving only the cutoff changes the cost by far
more (Table~\ref{tab:d-cost} and Fig.~\ref{fig:d-f1}). The validation-selected
F1 threshold is already $20\%$ above the optimum at $r=1$, and $75\%$ above
it at $r=20$, because maximising F1 assumes the two error types cost the
same. At $r=1$ the best model is also only $8.3\%$ cheaper than the trivial
policy of calling nobody, since predicting no for every record already
reaches $88.31\%$ accuracy; the model earns its cost once a missed
subscriber is about three times as expensive as a wasted call.

\begin{table}[tb]
\caption{Cost of the deployed cutoff versus the cheapest achievable cost,
per $1{,}000$ test clients. ``Deployed'' is the logistic fit with the
highest F1 (MLE) at its validation-selected threshold; ``best'' re-picks
the model \emph{and} the threshold for each ratio. The trivial policies
cost $r\cdot 116.9$ (call nobody) and $883.1$ (call everybody).}
\label{tab:d-cost}
\centering
\setlength{\tabcolsep}{4.2pt}
\begin{tabular}{@{}lrrrr@{}}
\hline
$r$ & Best cost & Deployed cost & Excess & Below ``nobody'' \\
\hline
1  & 107.2  & 128.9  & $+20.2\%$ & $8.3\%$  \\
5  & 383.0  & 393.2  & $+2.7\%$  & $34.5\%$ \\
10 & 631.8  & 723.5  & $+14.5\%$ & $46.0\%$ \\
15 & 752.2  & 1053.8 & $+40.1\%$ & $57.1\%$ \\
20 & 789.1  & 1384.2 & $+75.4\%$ & $66.3\%$ \\
\hline
\end{tabular}
\end{table}

\begin{figure}[tb]
\centering
\includegraphics[width=\textwidth]{s1_f1_vs_threshold.png}
\caption{F1 against the probability cutoff on the test split, one curve per
model; the marker is the validation-selected threshold actually deployed.
The logistic fits peak in a narrow band and collapse away from it, whereas
the three naive-Bayes models are flat between $0.2$ and $0.8$: their
``optimal'' cutoff is not identifiable, and consequently neither is any
prior-shift adjustment of it.}
\label{fig:d-f1}
\end{figure}

\paragraph{The base rate changes the metric, not the ranking.}
Reweighting the test set to positive rates between $2\%$ and $60\%$ leaves
ROC-AUC numerically constant (range $6.4\times10^{-14}$ over sixty
settings), while the best model's PR-AUC moves from $0.241$ at $5\%$ to
$0.784$ at $50\%$. Measured against the random baseline, which equals the
positive rate itself, the same model is $4.82$ times better than chance at
$5\%$ and only $1.57$ times better at $50\%$
(Fig.~\ref{fig:d-shift}). The rank order of the five models is unchanged at
every base rate. The practical reading is twofold: a bare
``PR-AUC $=0.404$'' is not interpretable without its $11.69\%$ base rate
and should be quoted as a multiple of the baseline, and the model is worth
deploying mainly where subscriptions are rare. Re-tuning the cutoff for a
shifted base rate is well behaved for the logistic fits, whose optimum
moves monotonically from $0.244$ at $5\%$ to $0.055$ at $50\%$, but is
ill-defined for the naive-Bayes models, whose optimum jumps between
$1.000$ and $0.0003$ across the same range.

% Optional: delete this figure and keep the numbers in the text if the
% section would otherwise exceed its page budget.
\begin{figure}[tb]
\centering
\includegraphics[width=\textwidth]{s3_prior_shift.png}
\caption{Prior-shift sensitivity. Left: ROC-AUC is invariant by
construction, because both of its axes are within-class rates. Right:
PR-AUC rises with the positive rate because precision mixes both classes;
the dashed line is the random baseline, which is the positive rate itself.
The dotted line marks the $11.69\%$ rate of this split.}
\label{fig:d-shift}
\end{figure}
